\documentclass[10pt,a4paper]{article}

\usepackage[utf8]{inputenc}
\usepackage[T1]{fontenc}
\usepackage[french]{babel}
\usepackage{amsmath,amssymb}
\usepackage{geometry}
\usepackage{enumitem}
\usepackage{multicol}
\usepackage{xcolor}
\usepackage{mdframed}
\usepackage{lmodern}
\usepackage{microtype}
\usepackage{fancyhdr}
\usepackage{titlesec}
\usepackage{parskip}

\geometry{top=1.8cm, bottom=2cm, left=1.8cm, right=1.8cm}

\pagestyle{fancy}
\fancyhf{}
\fancyhead[L]{\small\textit{Mathématiques, fiche de révision}}
\fancyhead[R]{\small\textit{Intégrales multiples}}
\fancyfoot[C]{\small\thepage}
\fancyfoot[R]{\tiny\textcopyright\ Rohan Fossé}
\renewcommand{\headrulewidth}{0.3pt}

\definecolor{bleu}{RGB}{30,60,120}
\definecolor{vert}{RGB}{0,110,60}
\definecolor{rouge}{RGB}{160,20,20}
\definecolor{gris}{gray}{0.93}
\definecolor{jauneclair}{RGB}{255,250,220}
\definecolor{vertclair}{RGB}{220,245,220}
\definecolor{bleutresleger}{RGB}{230,240,255}

\newmdenv[backgroundcolor=gris, linewidth=0pt,
  innerleftmargin=8pt, innerrightmargin=8pt,
  innertopmargin=5pt, innerbottommargin=5pt,
  skipabove=4pt, skipbelow=4pt]{grisbox}

\newmdenv[backgroundcolor=jauneclair, linecolor=rouge, linewidth=1pt,
  innerleftmargin=8pt, innerrightmargin=8pt,
  innertopmargin=5pt, innerbottommargin=5pt,
  skipabove=4pt, skipbelow=4pt]{alertbox}

\newmdenv[backgroundcolor=bleutresleger, linecolor=bleu, linewidth=1pt,
  innerleftmargin=8pt, innerrightmargin=8pt,
  innertopmargin=5pt, innerbottommargin=5pt,
  skipabove=4pt, skipbelow=4pt]{theobox}

\titleformat{\section}[block]{\normalsize\bfseries\color{bleu}}{}{0pt}{}[\vspace{-3pt}\rule{\linewidth}{0.5pt}]
\titlespacing{\section}{0pt}{10pt}{4pt}

\titleformat{\subsection}[runin]{\small\bfseries\color{vert}}{}{0pt}{}[\ :]
\titlespacing{\subsection}{0pt}{4pt}{0pt}

\begin{document}

\begin{center}
  {\Large\bfseries\color{bleu} Fiche de révision : Intégrales multiples}\\[3pt]
  {\small Intégrales doubles · Changements de variables · Intégrales triples · Green-Riemann}
\end{center}
\noindent\rule{\linewidth}{0.8pt}

\begin{multicols}{2}

\section{Intégrale double, théorème de Fubini}

\begin{theobox}
\textbf{Théorème de Fubini.} Sur un domaine $D=[a,b]\times[c,d]$ (rectangle) :
$$\iint_D f(x,y)\,dA = \int_a^b\!\left(\int_c^d f(x,y)\,dy\right)dx$$
$$= \int_c^d\!\left(\int_a^b f(x,y)\,dx\right)dy$$

\textbf{Cas séparable :} si $f(x,y)=g(x)\cdot h(y)$ sur un rectangle :
$$\iint_D f\,dA = \left(\int_a^b g(x)\,dx\right)\cdot\left(\int_c^d h(y)\,dy\right)$$
\end{theobox}

\section{Domaine non rectangulaire, inversion}

\begin{grisbox}
Sur un domaine général $D$, on intègre en identifiant les bornes variables.

\textbf{Inversion de l'ordre :} le domaine $\{0\le x\le 1,\; x\le y\le 1\}$ est aussi $\{0\le y\le 1,\; 0\le x\le y\}$.

$$\int_0^1\!\int_x^1 f(x,y)\,dy\,dx = \int_0^1\!\int_0^y f(x,y)\,dx\,dy$$

\textbf{Méthode :} dessiner le domaine $D$, puis relire les bornes dans l'autre sens.
\end{grisbox}

\section{Coordonnées polaires}

\begin{grisbox}
$$x = r\cos\theta,\quad y = r\sin\theta,\quad dA = r\,dr\,d\theta$$

\textbf{À utiliser quand :} domaine circulaire ou $\sqrt{x^2+y^2}$ apparaît dans $f$.

Sur le disque $x^2+y^2\le R^2$ : $r\in[0,R]$, $\theta\in[0,2\pi]$.
\end{grisbox}

\section{Formule de Green-Riemann}

\begin{theobox}
Si $D$ est un domaine fermé borné à bord $\partial D$ orienté positivement (sens trigonométrique) :
$$\boxed{\oint_{\partial D} P\,dx + Q\,dy = \iint_D \left(\frac{\partial Q}{\partial x}-\frac{\partial P}{\partial y}\right)dA}$$

\textbf{Application -- Calcul d'aire :} avec $P=-y$, $Q=x$ :
$$\mathrm{Aire}(D) = \frac{1}{2}\oint_{\partial D}(-y\,dx+x\,dy)$$
\end{theobox}

\section{Intégrales triples, coordonnées sphériques}

\begin{grisbox}
$$x=\rho\sin\varphi\cos\theta,\quad y=\rho\sin\varphi\sin\theta,\quad z=\rho\cos\varphi$$
$$dV = \rho^2\sin\varphi\;d\rho\;d\varphi\;d\theta$$

Bornes habituelles : $\rho\ge 0$, $\varphi\in[0,\pi]$, $\theta\in[0,2\pi]$.

\textbf{Coordonnées cylindriques :}
$$x=r\cos\theta,\ y=r\sin\theta,\ z=z,\quad dV=r\,dr\,d\theta\,dz$$
\end{grisbox}

\section{Applications}

\textbf{Centre de gravité} (densité $\rho=1$, aire $A=\iint_D dA$) :
$$\bar{x} = \frac{1}{A}\iint_D x\,dA \qquad \bar{y} = \frac{1}{A}\iint_D y\,dA$$

\section{Récapitulatif : quand utiliser quoi ?}

\begin{grisbox}
\renewcommand{\arraystretch}{1.3}
\begin{tabular}{@{}lp{5cm}@{}}
\textbf{Situation} & \textbf{Méthode} \\
\hline
Domaine rectangulaire & Fubini directement \\
$f=g(x)\cdot h(y)$ & Séparation de Fubini \\
Disque, anneau & Coordonnées polaires \\
Boule, sphère & Coordonnées sphériques \\
Cylindre & Coordonnées cylindriques \\
Calcul d'aire par contour & Green-Riemann \\
\end{tabular}
\end{grisbox}

\begin{alertbox}
\textbf{Jacobiens à ne pas oublier !}
\begin{itemize}[leftmargin=1em,itemsep=0pt]
  \item Polaires : $dA = r\,dr\,d\theta$ (facteur $r$)
  \item Sphériques : $dV = \rho^2\sin\varphi\,d\rho\,d\varphi\,d\theta$
  \item Cylindriques : $dV = r\,dr\,d\theta\,dz$
\end{itemize}
\end{alertbox}

\end{multicols}

\end{document}
